\section{Finite-rate certificates from waiting-time bins}\label{app:bins}
This appendix proves Proposition~\ref{prop:bins}. It states the variational inequality used, then the two certified computations: an elementary rational bound for fast reporters ($R=10^4,10^5$) and a validated enclosure of the finite-rate dynamics ($R=750,1000$). The scripts \texttt{verify\_lag\_bins\_exact.py} and \texttt{verify\_moderate\_speed.py} carry out the arithmetic exactly; \texttt{lag\_bins\_proof.md} in the code archive records the same argument with the implementation details.

\paragraph{The variational inequality.} Let $P$ be the forward law of an ordered pair of consecutive marks with the waiting time between them, and $Q$ the backward law with the two marks reversed and exchanged. In this example every directed mark and its counterpart have the same frequency in both experiments, so by the proof of Corollary~\ref{cor:pump} the estimator equals $K\KL(P\|Q)$, where $K$ is the activity of all observed jumps ($R/2$ for one edge, $R/2+H$ for all six). Choose positive bin weights $w_\ell$, assign $w_\ell$ to a pair of consecutive reporter marks whose rescaled wait $u=R\tau$ lies in bin $\ell$, and weight one to every pair containing a slow mark. This test is invariant under reversal of the pair. The variational formula for relative entropy and $\log z\le z-1$ give
\[
 K\KL(P\|Q)\ge K\{\mathbb E_P\log w-\mathbb E_Q(w-1)\}.
\]
Only pairs that start with a reporter event contribute, so both expectations may be taken conditional on a reporter start and multiplied by $R/2$. Under complete observation, a slow event occurring first sets the score to zero.

\paragraph{Fast reporters.} The bins have the eight finite thresholds $1/4,1/2,1,2,3,5,8,12$ and the nine weights $(90,94,101,114,132,132,103,77,56)/100$. As a comparison process keep the reporter enabled after the initial event, independently of the slow state. Given the starting phase $s$, its survival is $\exp[-R\int_0^{u/R}q(s+v)\dd v]$, whose logarithm differs from that of the frozen survival $e^{-uq(s)}$ by at most $\epsilon=\pi(16/17)u^2/R$. Multiplying enclosures of the frozen mean survival by $1-\epsilon$ and $1/(1-\epsilon)$ therefore encloses the true mean survival whenever $\epsilon<1$. Let $L(s)$ be the mean waiting time to this always-enabled reporter; the periodic backward equation $L'=RqL-1$ gives $\av{qL}=1/R$. The probability that a slow exit precedes the comparison reporter, averaged over the phase of a typical reporter event, is at most $2\|b\|_\infty\av{pqL}\le 2\|b\|_\infty(4/5)/R$, using $p\le 1/2+h\sqrt2/4<4/5$. Coupling the actual observation to the comparison process up to the first slow exit changes each expected score by at most this probability times the range of the score (which contains zero). The frozen mean survivals $\av{2pq\,e^{-uq}}$ at the thresholds are enclosed by the Taylor polynomials of degrees $111$ and $112$ of the exponential, whose phase averages follow from $\av{\cos^{2j}\theta}=\binom{2j}{j}/4^j$, vanishing odd moments and $\cos2\theta=2\cos^2\theta-1$. A step score has expectation $v_0+\sum_\ell(v_{\ell+1}-v_\ell)\Pr[U>u_\ell]$, and each survival enclosure enters at the endpoint dictated by the sign of its coefficient. Logarithms are bounded below by truncating the series of $2\operatorname{artanh}[(w-1)/(w+1)]$ and subtracting a geometric bound on the tail, and $\pi<22/7$ is used throughout. The result is a common lower bound on both estimators above $21.62$ at $R=10^4$ and above $366.44$ at $R=10^5$; with $\sigma\le88/35$ the ratio exceeds $145.74$.

\paragraph{Moderate rates.} Here the bins end at $1/4,1/2,3/4,1,3/2,2,3,4,6,10,16$ with weights $(90,94,99,104,111,121,132,136,121,90,67)/100$, and waits beyond $16$ receive weight one. In rescaled time the phase moves with speed $\delta=2\pi/R$. On periodic pairs of functions (twins, singleton) the killed backward generator is
\[
 A f = \delta\partial_\theta f+
 \begin{pmatrix}-q-b/R & b/R\\ a/R&-a/R\end{pmatrix} f ,
\]
in which a reporter jump kills the process; for complete observation a slow exit also kills, and $A=\delta\partial_\theta-q-b/R$ acts on the twins alone. With $g=(q,0)^{\top}$ (or $q$), the probability of a reporter event by rescaled time $u$, starting from the twins at phase $\theta$, is the twin component of $F(u)=\int_0^u e^{vA}g\dd v$, and $e^{vA}$ contracts the supremum norm. Put $U=16$, $V_0=Ug$ and $V_{n+1}=U A V_n/(n+2)$, so that $V_n=U^{n+1}A^n g/(n+1)!$. Taylor's formula with integral remainder and contraction give, for $0\le u\le U$,
\[
 \Big\|F(u)-\sum_{n=0}^{N-1}(u/U)^{n+1}V_n\Big\|_\infty
 \le (u/U)^{N+1}\|V_N\|_\infty .
\]
Each $V_n$ is a trigonometric polynomial, so the recurrence needs only exact differentiation and multiplication of Fourier series. The script stores real and imaginary coefficient intervals as integers over the common denominator $2^{224}$, rounds every operation outward, encloses $\pi$ by Machin's formula, bounds the supremum norm by the sum of absolute coefficients, and stops when the enclosure of the remainder is below $10^{-12}$. Projecting on the event-phase density $2pq$ gives enclosures of the bin masses, and the variational inequality completes the bound. No floating-point operation or ordinary differential equation solver enters. The certified bounds are $2.746942$ (one edge) and $2.755619$ (all jumps) at $R=750$, and $3.654098$ and $3.661074$ at $R=1000$, each above $4\pi/5<88/35$.
